% ============================================================
% CAPÍTULO 6: CONCLUSÃO
% ============================================================

\chapter{Conclusão}
\label{cap:conclusao}

\section{Retomada do Problema e das Hipóteses}
\label{sec:retomada-hipoteses}

Esta monografia partiu do seguinte problema de pesquisa: \textit{de que forma a Encíclica \MH{} pode contribuir para a proteção jurídica da personalidade humana diante dos desafios impostos pelo desenvolvimento e difusão da inteligência artificial?}

Para responder a essa questão, foram formuladas quatro hipóteses de trabalho:

\begin{itemize}
    \item \textbf{H1}: A \MH{} oferece uma fundamentação antropológica robusta para a proteção da personalidade, ancorada no personalismo cristão e em uma concepção relacional e transcendental da pessoa.
    \item \textbf{H2}: A IA apresenta riscos específicos e qualitativamente novos à personalidade humana, que transcendem as categorias tradicionais de proteção jurídica.
    \item \textbf{H3}: O ordenamento jurídico brasileiro dispõe de instrumentos relevantes para a proteção da personalidade, mas estes apresentam lacunas e insuficiências diante dos desafios específicos da IA.
    \item \textbf{H4}: A \MH{} oferece contribuições normativas que podem orientar o aperfeiçoamento da regulação da IA no Brasil, especialmente nos planos principiológico, crítico e institucional.
\end{itemize}

O percurso investigativo realizado nos capítulos precedentes permite agora avaliar em que medida essas hipóteses foram confirmadas.

\section{Verificação das Hipóteses}
\label{sec:verificacao-hipoteses}

\subsection{H1: A Fundamentação Antropológica da \MH{}}
\label{subsec:verificacao-h1}

A análise desenvolvida no capítulo segundo confirmou integralmente a primeira hipótese. A \MH{} oferece, de fato, uma fundamentação antropológica robusta para a proteção da personalidade, articulada em três dimensões: ontológica (a pessoa como ser dotado de dignidade intrínseca), relacional (a pessoa como ser constituído nas e pelas relações com os outros) e transcendental (a pessoa como ser aberto à transcendência).

O conceito de ``irredutibilidade da pessoa aos seus dados'', que a Encíclica desenvolve em diversos parágrafos (§3, §101, §186, §232), mostrou-se particularmente relevante para o debate contemporâneo sobre proteção de dados e regulação da IA. Esse conceito oferece um fundamento antropológico para a tese de que a pessoa não pode ser reduzida a seus traços digitais, a seus perfis algorítmicos ou a suas classificações preditivas.

Os cinco princípios éticos identificados (dignidade, subsidiariedade, solidariedade, transparência e responsabilidade) constituem um quadro normativo coerente e aplicável à regulação da IA.

\subsection{H2: Os Riscos Específicos da IA}
\label{subsec:verificacao-h2}

A análise desenvolvida no capítulo terceiro confirmou a segunda hipótese, demonstrando que a IA apresenta riscos qualitativamente novos à personalidade humana. A identificação de quatro dimensões de risco (a automação decisória e o direito à explicação, a discriminação algorítmica, o capitalismo de vigilância e a identidade pessoal em questão, e os desafios do transumanismo) evidenciou que a proteção da personalidade na era da IA não pode ser adequadamente enfrentada com os instrumentos tradicionais do direito civil e constitucional.

Particularmente significativa foi a constatação de que a opacidade dos sistemas algorítmicos desafia a própria possibilidade de exercício de direitos, na medida em que a pessoa afetada por uma decisão automatizada frequentemente não dispõe dos meios para sequer identificar a existência da violação.

\subsection{H3: As Insuficiências do Ordenamento Brasileiro}
\label{subsec:verificacao-h3}

A análise do capítulo quarto confirmou a terceira hipótese, demonstrando que, embora o ordenamento jurídico brasileiro disponha de instrumentos relevantes, estes apresentam lacunas significativas. A LGPD, apesar de seu caráter inovador, carece de instrumentos específicos para o controle da discriminação algorítmica e para a garantia efetiva do direito à explicação. O Marco Civil da Internet não oferece resposta adequada para a responsabilidade civil por conteúdos gerados por IA. O PL 2.338/2023, embora promissor, ainda está em tramitação e seu texto final é incerto.

A principal lacuna identificada foi a ausência, nos instrumentos jurídicos brasileiros, de uma fundamentação antropológica explícita que oriente a interpretação e aplicação das normas de proteção da pessoa diante dos riscos algorítmicos.

\subsection{H4: As Contribuições Regulatórias da \MH{}}
\label{subsec:verificacao-h4}

O capítulo quinto confirmou a quarta hipótese, demonstrando que a \MH{} oferece contribuições normativas significativas para o aperfeiçoamento da regulação da IA. No plano crítico, fornece instrumentos conceituais para questionar o paradigma tecnocrático e a tese da neutralidade tecnológica. No plano principiológico, oferece fundamentos antropológicos robustos para dignidade, transparência, subsidiariedade e solidariedade como critérios regulatórios. No plano institucional, aponta direções para o aperfeiçoamento da LGPD e para o enriquecimento do PL 2.338/2023.

A tese central, de que a \MH{} pode fundamentar uma regulação da IA que vá além da mera gestão de riscos e se constitua como verdadeira proteção da pessoa humana, foi confirmada pela análise.

\section{Resposta ao Problema de Pesquisa}
\label{sec:resposta-problema}

À luz das evidências e análises apresentadas, é possível oferecer uma resposta ao problema de pesquisa proposto.

A \MH{} contribui para a proteção jurídica da personalidade humana diante dos desafios da IA em três planos fundamentais. Em \textbf{primeiro lugar}, oferece uma fundamentação antropológica para os direitos da personalidade que supera as concepções meramente formalistas ou procedimentais predominantes no debate regulatório contemporâneo. A Encíclica recorda que a proteção da pessoa não é uma questão técnica a ser resolvida por mecanismos de compliance, mas uma exigência ética fundamental que deve orientar todo o desenvolvimento tecnológico.

Em \textbf{segundo lugar}, a \MH{} formula princípios éticos (dignidade, transparência, subsidiariedade, solidariedade, responsabilidade) que podem operar como critérios substantivos para a avaliação e o aperfeiçoamento dos instrumentos jurídicos de proteção da personalidade. Esses princípios não substituem o trabalho técnico da regulação, mas conferem-lhe uma orientação axiológica clara e consistente.

Em \textbf{terceiro lugar}, a \MH{} oferece uma crítica ao paradigma tecnocrático que permite compreender os riscos da IA não como fenômenos isolados e corrigíveis, mas como manifestações de uma visão de mundo que tende a reduzir a pessoa a seus aspectos mensuráveis e funcionais. Essa crítica fundamenta a exigência de que a regulação da IA não se limite a controlar abusos, mas promova ativamente um desenvolvimento tecnológico centrado na pessoa.

\section{Limitações da Pesquisa}
\label{sec:limitacoes}

Este trabalho apresenta limitações que devem ser explicitadas para a adequada compreensão de seus resultados.

Em \textbf{primeiro lugar}, a pesquisa limitou-se ao exame do ordenamento jurídico brasileiro, não abrangendo o direito comparado nem o direito internacional, salvo referências pontuais ao GDPR europeu. Uma investigação mais ampla poderia identificar contribuições adicionais da \MH{} para outros sistemas jurídicos.

Em \textbf{segundo lugar}, a análise da regulação da IA concentrou-se em instrumentos normativos já existentes ou em tramitação (LGPD, MCI, PL 2.338/2023), sem aprofundar o exame de normas técnicas, autorregulatórias ou de soft law que também compõem o ecossistema regulatório.

Em \textbf{terceiro lugar}, a pesquisa adotou uma abordagem predominantemente teórica e analítica, sem realizar estudos empíricos sobre a aplicação concreta dos instrumentos normativos examinados. Estudos de caso sobre decisões judiciais ou administrativas envolvendo IA e direitos da personalidade poderiam oferecer uma perspectiva complementar relevante.

Em \textbf{quarto lugar}, a \MH{} é um documento recente (publicado em 25 de maio de 2026), e a literatura secundária sobre ela ainda é incipiente. À medida que a recepção acadêmica e eclesial da Encíclica se desenvolva, novas interpretações e aplicações poderão enriquecer o debate aqui iniciado.

\section{Agendas de Pesquisa Futura}
\label{sec:agendas-futuras}

Os resultados deste trabalho abrem diversas possibilidades de investigação futura, que podem ser organizadas em quatro eixos.

O \textbf{primeiro eixo} diz respeito ao aprofundamento teórico dos conceitos aqui explorados. A irredutibilidade da pessoa aos dados, o paradigma tecnocrático e os princípios éticos da \MH{} merecem desenvolvimento sistemático em diálogo com a filosofia do direito, a teoria dos direitos fundamentais e a ética aplicada.

O \textbf{segundo eixo} concerne à análise empírica da aplicação dos instrumentos jurídicos examinados. Estudos de caso sobre decisões judiciais envolvendo discriminação algorítmica, direito à explicação e responsabilidade civil por danos causados por sistemas de IA poderiam oferecer uma perspectiva complementar à análise teórica aqui empreendida.

O \textbf{terceiro eixo} refere-se ao direito comparado. A investigação sobre como outros ordenamentos jurídicos (especialmente o europeu (AI Act), o norte-americano e o chinês) têm enfrentado os desafios da IA à personalidade humana poderia oferecer subsídios valiosos para o aperfeiçoamento da regulação brasileira.

O \textbf{quarto eixo} diz respeito ao desenvolvimento tecnológico. À medida que novas aplicações de IA surjam e que os sistemas existentes se tornem mais sofisticados, novas questões de proteção da personalidade certamente emergirão, exigindo permanente atualização do marco regulatório e da reflexão ética que o fundamenta.

\section{Considerações Finais}
\label{sec:consideracoes-finais}

A inteligência artificial não é, em si mesma, uma ameaça à personalidade humana. Como recorda a \MH{} (§89), a tecnologia é uma expressão da criatividade humana e pode servir ao desenvolvimento integral da pessoa quando orientada por valores éticos adequados. O problema não está na tecnologia em si, mas na visão de mundo que a informa e nos usos que dela se fazem.

A \MH{} oferece uma contribuição singular para este debate não porque propõe soluções técnicas ou regras específicas, mas porque recorda o fundamento último de toda proteção jurídica da pessoa: a dignidade que lhe é inerente, anterior a qualquer reconhecimento estatal, independente de qualquer classificação algorítmica.

O desafio que se coloca aos juristas, legisladores e operadores do direito é o de traduzir essa fundamentação antropológica em instrumentos normativos eficazes, capazes de proteger a pessoa humana na era da inteligência artificial. Este trabalho espera ter oferecido uma contribuição para esse esforço coletivo, cuja urgência e importância só tendem a crescer à medida que a IA se integra cada vez mais profundamente em todos os âmbitos da vida social.
